% ============================================================
% Section IV: Implementation Details
% ============================================================

\section{Implementation Details}
\label{sec:implementation}

\subsection{Technology Stack}
Table~\ref{tab:tech_stack} summarizes the technology stack employed in AgriSense.

\begin{table}[htbp]
\caption{Technology Stack of AgriSense}
\label{tab:tech_stack}
\centering
\begin{tabular}{@{}ll@{}}
\toprule
\textbf{Component} & \textbf{Technology} \\
\midrule
Backend Framework & FastAPI (Python 3.9+) \\
ML Library & scikit-learn (Random Forest) \\
Data Handling & pandas, NumPy \\
Model Serialization & pickle \\
Generative AI & Google Gemini 2.5 Flash \\
Frontend Framework & React 18 + Vite \\
HTTP Client & Axios \\
Visualization & Recharts (Pie, Bar Charts) \\
UI Icons & Lucide React \\
State Management & React Hooks (useState, useEffect) \\
Routing & React Router DOM \\
API Communication & RESTful JSON over HTTP \\
Cross-Origin Support & FastAPI CORSMiddleware \\
\bottomrule
\end{tabular}
\end{table}

\subsection{Backend Architecture}
The backend is structured following a \textbf{Service-Oriented Architecture (SOA)} pattern with clear separation of concerns:

\begin{itemize}
    \item \textbf{Routes Layer}: Defines API endpoints using FastAPI's \texttt{APIRouter}. Each module has a dedicated route file (\texttt{recommend.py}, \texttt{yield.py}, \texttt{profit.py}, \texttt{risk.py}, \texttt{compare.py}) that handles HTTP request validation via Pydantic \texttt{BaseModel} schemas.
    \item \textbf{Services Layer}: Encapsulates business logic. Each service class (e.g., \texttt{RecommendationService}, \texttt{YieldService}) manages model loading, preprocessing, and inference independently.
    \item \textbf{Models Layer}: Stores serialized ML model artifacts (\texttt{.pkl} files) generated by offline training scripts.
    \item \textbf{AI Integration}: A dedicated \texttt{/api/insight} endpoint receives contextual prompts and interfaces with the Gemini API via the \texttt{google-generativeai} Python SDK.
\end{itemize}

\subsection{Frontend Architecture}
The frontend is built as a Single-Page Application (SPA) using React 18 with Vite as the build tool. Key architectural decisions include:

\begin{itemize}
    \item \textbf{Component-Based Design}: Reusable components such as \texttt{Navbar} and \texttt{AIInsightButton} promote code reuse and consistency across pages.
    \item \textbf{Page-Based Routing}: React Router DOM manages navigation between modules (Home, CropRecommendation, YieldPrediction, ProfitAnalysis, RiskAssessment, CompareCrops).
    \item \textbf{Interactive Visualizations}: Recharts library provides responsive, interactive charts---Pie Charts for crop probability distributions and Bar Charts for yield and profit comparisons.
    \item \textbf{Authentication}: A localStorage-based session management mechanism gates access to the dashboard.
    \item \textbf{Glassmorphism UI}: A modern, dark-themed UI employing glassmorphism design principles with CSS custom properties for theming.
\end{itemize}

\subsection{API Design}
Table~\ref{tab:api_endpoints} documents the RESTful API endpoints exposed by the backend.

\begin{table}[htbp]
\caption{AgriSense RESTful API Endpoints}
\label{tab:api_endpoints}
\centering
\begin{tabular}{@{}lll@{}}
\toprule
\textbf{Endpoint} & \textbf{Method} & \textbf{Description} \\
\midrule
\texttt{/api/recommend} & POST & Crop recommendation \\
\texttt{/api/yield} & POST & Yield prediction \\
\texttt{/api/profit} & POST & Profit calculation \\
\texttt{/api/risk} & POST & Risk assessment \\
\texttt{/api/compare} & POST & Crop comparison \\
\texttt{/api/insight} & POST & AI insight generation \\
\bottomrule
\end{tabular}
\end{table}

\subsection{Training Pipeline}
Two offline training scripts are provided:

\begin{itemize}
    \item \texttt{train\_classifier.py}: Loads the crop recommendation dataset, encodes labels using \texttt{LabelEncoder}, trains a Random Forest Classifier with 100 estimators, and serializes both the model and encoder.
    \item \texttt{train\_regressor.py}: Loads the crop yield dataset, constructs a \texttt{Pipeline} with \texttt{ColumnTransformer} (One-Hot Encoding) and Random Forest Regressor (50 estimators), and serializes the entire pipeline.
\end{itemize}

Both scripts use a fixed random seed (\texttt{random\_state=42}) for reproducibility.
